% Part: proof-theory
% Chapter: natural-deduction
% Section: quantifiers

\documentclass[../../../include/open-logic-section]{subfiles}

\begin{document}

\olfileid{pt}{nat}{qua}

\olsection{\usetoken{P}{proof} Reguler lan Substitusi}

Aturan kuantor njalari sawatara keruwetan amarga
``syarat variabel eigen'' kang lumaku kanggo aturan
\Elim{\lexists} lan~\Intro{\lforall}.
Akèh-akèhé keruwetan iki bisa diatasi kanthi njamin
yèn !!{constant}~$c$ ing inferènsi mau ora tau dienggo
manèh. Kita tetepaké definisi iki.

\begin{defn}
!!{proof}~$\delta$ ing deduksi alami diarani
\emph{reguler} yèn
\begin{enumerate}
  \item ora ana variabel eigen~$a$ kang dienggo ing
  luwih saka siji inferènsi \Elim{\lexists}
  utawa~\Intro{\lforall}, lan
  \item yèn $a$ variabel eigen ing inferènsi
  \Elim{\lexists} utawa~\Intro{\lforall},
  variabel iku mung ana ing sub-!!{proof} langsung
  kang nuju premis cilik, utawa premis tunggal,
  miturut urutané.
\end{enumerate}
\end{defn}

\begin{lem}\ollabel{lem:subst-regular}
Yèn $\delta$ reguler lan pungkasané ana ing~$!C$,
$c$ iku !!a{constant} kang ora dadi variabel eigen
ing~$\delta$, sarta $t$~suku kang ora ngemot variabel
eigen saka~$\delta$, banjur
$\Subst{\delta}{t}{c}$, kang diolehké saka~$\delta$
kanthi ngganti saben kedadéan~$c$ nganggo~$t$,
iku !!{proof} reguler. !!{formula} pungkasan saka
$\Subst{\delta}{t}{c}$ yaiku $\Subst{!C}{t}{c}$.
Yèn $!A^x$ asumsi kabuka ing~$\delta$, banjur
$\Subst{!A}{t}{c}^x$ asumsi kabuka ing
~$\Subst{\delta}{t}{c}$.
\end{lem}

\begin{proof}
  Buktiné nganggo induksi tumrap $\pheight{\delta}$.
  Yèn $\pheight{\delta} = 0$, isiné mung asumsi~$!A^x$.
  Banjur $\Subst{\delta}{t}{c}$ padha karo
  ~$\Subst{!A}{t}{c}^x$.

  Yèn $\pheight{\delta} > 0$, kita bedakaké manut
  inferènsi pungkasan ing~$\delta$. Kahanan aturan
  konektif proposisional gampang lan dipasrahaké
  dadi latihan. Kita ngrembug nalika $\delta$
  pungkasané nganggo inferènsi kuantor.
  \begin{enumerate}
    \item Yèn inferènsi pungkasan yaiku $\Intro{\lforall}$,
    $\delta$ wujudé
    \[
    \AxiomC{}
    \RightLabel{$\delta'$}
    \DeduceC{$!A(a,c)$}
    \RightLabel{\Intro{\lforall}}
    \UnaryInfC{$\lforall[x][!A(x,c)]$}
    \DisplayProof
    \]
    Miturut hipotesis induksi, $\Subst{\delta'}{t}{c}$
    iku !!{proof} reguler kanggo $!A(a,t)$.
    Amarga $a$ variabel eigen ing~$\delta$,
    $a$~ora ana ing~$t$. Mula inferènsi pungkasan ing
    $\Subst{\delta}{t}{c}$,
    \[
    \AxiomC{}
    \RightLabel{$\Subst{\delta'}{t}{c}$}
    \DeduceC{$!A(a,t)$}
    \RightLabel{\Intro{\lforall}}
    \UnaryInfC{$\lforall[x][!A(x,t)]$}
    \DisplayProof
    \]
    iku bener: amarga $t$ ora ngemot~$a$, konklusi
    lan kabèh asumsi kabuka uga ora ngemot~$a$.
    \item Yèn inferènsi pungkasan yaiku
    $\Elim{\lforall}$, $\delta$ wujudé
    \[
    \AxiomC{}
    \RightLabel{$\delta'$}
    \DeduceC{$\lforall[x][!A(x,c)]$}
    \RightLabel{\Elim{\lforall}}
    \UnaryInfC{$!A(s(c),c)$}
    \DisplayProof
    \]
    Miturut hipotesis induksi,
    $\Subst{\delta'}{t}{c}$ iku !!{proof} reguler kanggo
    $\lforall[x][!A(x,t)]$. (Suku~$s(c)$ ing konklusi
    bisa ngemot~$c$, nanging ora kudu.)
    Inferènsi pungkasan ing $\Subst{\delta}{t}{c}$,
    \[
    \AxiomC{}
    \RightLabel{$\Subst{\delta'}{t}{c}$}
    \DeduceC{$\lforall[x][!A(x,t)]$}
    \RightLabel{\Elim{\lforall}}
    \UnaryInfC{$!A(s(t),t)$}
    \DisplayProof
    \]
    iku bener, lan
    $!A(s(t),t) \ident \Subst{!A(s(c),c)}{t}{c}$.
    Kahanan \Intro{\lexists} padha polané.
  \item Yèn inferènsi pungkasan yaiku
    $\Elim{\lexists}$, $\delta$ wujudé
    \[
    \AxiomC{}
    \RightLabel{$\delta_1'$}
    \DeduceC{$\lexists[x][!A(x,c)]$}
    \AxiomC{$\Discharge{!A(a,c)}{x}$}
    \RightLabel{$\delta_2'$}
    \DeduceC{$!C$}
    \DischargeRule{\Elim{\lexists}}{x}
    \BinaryInfC{$!C$}
    \DisplayProof
    \]
    Miturut hipotesis induksi, $\Subst{\delta_1'}{t}{c}$
    lan $\Subst{\delta_2'}{t}{c}$ iku !!{proof}s reguler
    kanggo $\lexists[x][!A(x,t)]$ lan
    $\Subst{!C}{t}{c}$ saka asumsi kabuka~$!A(a,t)$.
    Amarga $a$ variabel eigen ing~$\delta$,
    $a$~ora ana ing~$t$. Mula inferènsi pungkasan ing
    $\Subst{\delta}{t}{c}$,
    \[
    \AxiomC{}
    \RightLabel{$\Subst{\delta_1'}{t}{c}$}
    \DeduceC{$\lexists[x][!A(x,t)]$}
    \AxiomC{$\Discharge{!A(a,t)}{x}$}
    \RightLabel{$\Subst{\delta_2'}{t}{c}$}
    \DeduceC{$\Subst{!C}{t}{c}$}
    \DischargeRule{\Elim{\lexists}}{x}
    \BinaryInfC{$\Subst{!C}{t}{c}$}
    \DisplayProof
    \]
    iku bener: $\Subst{!A(x,t)}{a}{x} \ident
    \Subst{!A(a,c)}{t}{c}$, lan konklusi
    ~$\Subst{!C}{t}{c}$ sarta kabèh asumsi kabuka
    ora ngemot~$a$.
  \end{enumerate}
\end{proof}

\begin{prop}\ollabel{prop:regular}
Yèn $\delta$ iku !!a{proof} ing \Log{N1c} utawa
\Log{N1i}, banjur ana !!{proof} reguler~$\delta'$
kanthi struktur kang padha (mliginé !!{height} kang
padha), kang béda saka $\delta$ mung amarga ngganti
variabel eigen.
\end{prop}

\begin{proof}
Kanggo bukti iki waé, sebut inferènsi variabel eigen
ing~$\delta$ \emph{resik} yèn variabel eigené ora dadi
variabel eigen inferènsi liya lan ora ana ing pérangan
liya saka~$\delta$ saliyané sub-!!{proof} langsung
saka inferènsi iku; yèn ora, sebut \emph{reged}.
Upama $n$ cacah inferènsi variabel eigen reged
ing~$\delta$. Kita tuduhaké kanthi induksi tumrap~$n$
yèn $\delta$~bisa diowahi dadi !!{proof} reguler~$\delta'$
kang dikarepaké. Yèn $n=0$, kabèh inferènsi variabel
eigen ing~$\delta$ resik, mula $\delta$ wis reguler
lan ora ana kang perlu dibuktèkaké. Yèn ora, pilihen
inferènsi variabel eigen reged kang paling dhuwur,
upamané \Intro{\lforall}. Sub-!!{proof}~$\delta_1$
ing~$\delta$ kang pungkasané ana ing inferènsi iku
wujudé
    \[
    \AxiomC{}
    \RightLabel{$\delta_2$}
    \DeduceC{$!A(c)$}
    \RightLabel{\Intro{\lforall}}
    \UnaryInfC{$\lforall[x][!A(x)]$}
    \DisplayProof
    \]
Amarga $c$ variabel eigen, ora ana ing
~$\lforall[x][!A(x)]$ utawa asumsi kabuka.
Bukti $\delta_2$ reguler, amarga saben inferènsi
variabel eigen ing kono resik. Upama $a$ iku
!!a{constant} kang ora ana ing~$\delta$.
Miturut \olref{lem:subst-regular},
$\Subst{\delta_2}{a}{c}$ iku bukti reguler kanggo
$!A(a)$. Miturut syarat variabel eigen, ora ana
asumsi kabuka ing~$\delta_2$ kang ngemot~$c$,
mula ora ana asumsi kabuka ing
~$\Subst{\delta_2}{a}{c}$ kang ngemot~$a$.
Dadi kita bisa ngetrapaké~$\Intro{\lforall}$.
Yèn $\delta_1$ ing~$\delta$ diganti bukti $\delta_1'$,
    \[
    \AxiomC{}
    \RightLabel{$\Subst{\delta_2}{a}{c}$}
    \DeduceC{$!A(a)$}
    \RightLabel{\Intro{\lforall}}
    \UnaryInfC{$\lforall[x][!A(x)]$}
    \DisplayProof
    \]
kita wis ngganti inferènsi variabel eigen reged dadi
resik, lan bédané mung variabel eigen $c$ diganti
nganggo~$a$. Bukti asilé nduwèni paling akèh $n-1$
inferènsi variabel eigen reged, mula asilé ndhèrèk saka
hipotesis induksi.
\end{proof}

\begin{prob}
Jlentrehna kahanan ing bukti \olref{prop:regular}
nalika inferènsi variabel eigen reged kang paling
dhuwur yaiku~\Elim{\lexists}.
\end{prob}

Kita ora bakal kanthi eksplisit ngrujuk proposisi kang
lagi dibuktèkaké, nanging bakal nganggep kabèh
!!{proof}s kang dirembug reguler---yèn durung,
variabel eigen mesthi bisa diganti supaya reguler.

\Olref{lem:subst-regular} lan \olref{prop:regular}
uga bener kanggo \Log{N2c} lan~\Log{N2i}.
Buktiné dipasrahaké dadi latihan.

\end{document}
